% The builder separates the blocks below by their unique subsection headings.
\subsection{各向异性状态方程：长电子层出现，后期总率仍相近}\label{sec:synthesis-anisotropy}
\citet{F611}在中等导引场的force-free和Harris片中比较全PIC、各向异性电子混合模型和等温电子混合模型。最直接的新现象是：前两者出现跨数个$d_i$的电子电流层，等温例没有同样长的层；Harris片X线附近的面外磁场还增强到背景的约1.5--2倍。图6的沿片剖面为双峰，不能把所有增强都画成X点正中心的单峰。

这补足了前面Sullivan层长研究与压力闭合研究之间的联系：层长不仅会随耗散改变，也可能由电子各向异性支撑。force-free例的$F_e=\mu_0(p_{e\parallel}-p_{e\perp})/B^2$峰值约0.7；Harris长层处约大于0.5。但图3、7的后期率在三种模型中仍接近，说明\textbf{复现总率不等于复现长层及Hall场}。早期时钟并不一致，Harris比较还用了不同种子并作时间平移。

这一结果与\citet{Cassak2015}不能合并成“任意电子各向异性均能独自产生快重联”。两者的状态方程不同；F611明确指出，所用俘获插值闭合在该参数区主要改变结构，并不自行提供无Hall快支。其Harris实现还要沿二维磁通标签追踪上游值，并修改闭合的参考密度以避免中心非物理各向异性。适用于无导场和复杂三维几何的可迁移闭合仍未建立。第一章第\ref{sec:synthesis-closure-model}节给出状态方程、极限和曲率电流的逐步推导。

\subsection{有限长X线：Hall磁通输运产生受抑区及整体率下降}\label{sec:synthesis-finite}
\citet{F659}用三维PIC把薄片夹在两端厚片之间，阻止X线无限扩展。其薄片长度为31、12、8.4和$4d_i$；短于约$10d_i$时，率和离子喷流显著下降。长例也不是处处以同一速率重联：电子漂移侧形成活跃区，离子漂移侧保留较厚、较弱的受抑区；次级撕裂优先出现在前者。

作者逐项检查欧姆定律后，并未把现象简单归于一个新增三维压力梯度项。活跃区破冻结仍主要由非对角电子压力散度支撑；关键的几何反馈是电子沿电流反方向搬运重联法向场$B_z$，使另一侧缺少驱动出流和继续变薄所需的磁张力。将约$10\Omega_i^{-1}$的变薄时间乘该处量级为$\VA$的电子漂移，得到约$10d_i$的受抑长度。其图8中局地非理想电场约0.12，沿整条线积分并除以薄片长度却约0.068；两者同时成立。

\citet{F792}沿同一研究路线补做八个长度算例，以$\delta<2d_{e,\mathrm{local}}$定义活跃区。按活跃长度归一后，有活跃区的率回到$O(0.1)$；按整条X线长度归一则下降。进一步的动量预算显示，$J_y$和被搬运的$B_z$在第三方向错位，降低$J_yB_z$；因此喷流减弱不只是把二维出流转向第三方向。第一章第\ref{sec:synthesis-length-model}节重建长度比例与梯形重叠模型。

这两篇是\textbf{动理学对照和同一连续研究序列}。2019例质量比75，2020例25，没有完成共同的真实质量比收敛；图5标度使用本组模拟拟合的受抑长度、过渡宽度和速度幅度。最短两例没有活跃区，却仍出现初扰动引发的瞬态喷流，零输出近似解释不了；作者也明确没有严格建立最大喷流与总率的唯一关系。弱导场、不同$\beta$及多斜模是否保留相同阈值，仍需检验。

\subsection{Hall电场与离子加热：轨道证据应怎样读}\label{sec:synthesis-heating}
\citet{L272}把电子正则涡量描述、给定场试粒子与PIC图像连接起来，提出强面内电场梯度使离子失去规则回旋运动的机制。所示PIC中，梯度指标较大的区域与高离子温度区大致对应；导引场由0增至0.5时，强效应的范围和程度减小。它提供了比“释放了磁能所以会热”更具体的空间预测，但图6、7来自其2019前序工作，不能计为一套独立复现。

认真读方法还需保留三个区别。第一，试粒子图1、2把粒子\emph{能量分布的标准差}记作TE，它不是通常由中心速度二阶矩定义的温度。第二，倒抛物势的线性运动方程本来就有指数发散解；仅凭这个例子不能独立建立复杂混沌或不可逆熵增。第三，论文中由准中性采用不可压缩，以及把给定几何下的阈值称作普遍热化条件，都需要额外假设。第一章第\ref{sec:synthesis-heating-model}节给出带符号的轨道方程、倒抛物势系数复核和能量分解。

这些限制没有取消其PIC中的空间关联，却改变了结论强度：它支持一条可检验的失磁化与增能途径，尚不足以单独解释所有导场、三维和高$\beta$重联中的热化。与\citet{E512}的相混合—弱碰撞链相比，有限矩模型还缺乏对速度空间细化及熵产的直接约束；与第\ref{sec:synthesis-mrx}节MRX原文比较，也不能把后续解释倒写为实验当时已唯一证实的机制。

\subsection{嵌套四极：多个X点并不自动产生同样的图样}\label{sec:synthesis-nested}
\citet{L213}在EMHD中显示，强主X点的出流与弱次级X点的入流相互作用，可以形成一个不直接围绕单个X点的新外四极，包围原来的扩展四极。特别值得保留的对照是：三个同强度重联点并未产生相同嵌套结构。因此，多X点只是条件之一，\textbf{相对强度和相互作用顺序}也决定图样。

其与Cluster的比较依赖虚拟轨道：依次核对三个磁分量的过零、Hall峰以及法向电场的投影，而不是直接测到一张完整二维四极图。模拟片半厚约$7d_e$、观测约$3$--$5d_e$；法向场仅约1 nT，坐标方向不确定度可达20度，单航天器又不能唯一分离时间和空间变化。报告因此把它列为与EMHD解释相容的观测，保留替代结构和坐标误差。

这也给下一小节的尘埃嵌套四极提供一个必要对照：两者可以形似，但前者来自电子尺度多X点相互作用，后者含更慢的第三物种响应。要判别机制，应比较内外尺度比、建立时间和分物种电流，不能从“嵌套”二字直接识别模型。

\subsection{MRX先测到加速与下游升温，热化机制仍待区分}\label{sec:synthesis-mrx}
\citet{Yoo2013}直接测得MRX的面内电势井在下游变深、变宽，面内电场约500--800 V/m，离子跨越分离线后可加速至约$0.5\VA$。沿下游方向到约9 cm后，流速开始趋于饱和而离子温度上升；温度横向剖面比喷流更窄。图2电势、图3离子流与电功、图4光谱和温度必须连读，才构成“先加速、再下游升温”的测量链。

这一实验约束了能量进入离子的路径，却没有唯一选出随机加热。原文明确保留再磁化、下游碰撞、波粒作用、势阱反弹及装置边界等解释；下游离子均程降至2--5 cm，也不支持把全区域一概称为严格无碰撞。电子压力梯度在电势重建中不可忽略，不能把测得的整个面内电场都机械等同于单独的$\J\times\B/(en)$。

\citet{L272}后来用电场梯度估计解释同类实验，是对这些事实提出的机制候选。与本报告已核读的非对称能量收支实验\citep{Yamada2018}相连后，研究问题变为：不同闭合能否同时预测电势、喷流饱和位置、温升位置和分物种能流？仅使一个电场阈值超过一，还不足以完成这四项检验。
